%% Question bank chapter file - generated by bank_tools.py build (do not edit while the bank is still being built:
%% it is regenerated from the per-test data in Question Bank\data). Every question is one bankq block:
%%   \begin{bankq}[id=..., source={...}, paper=.., qno=.., type=..., subtopic={...}, answer={...}, answer_src=key|solved,
%%                 working={...}, difficulty=1-10, diff_parts={K S I C F}, diff_note={...},
%%                 group=<passage id or empty>, also_in={duplicate copies in other tests}]
%% difficulty = min(10, K+S+I+C+F-1): K concepts 1-3, S steps 1-4, I insight 0-2, C computation 0-2, F format 0-1
%%   ...question LaTeX, options, figures (\bfig / \bfigi -> img/)...
%%   \end{bankq}
%% Passages shared by several questions are bankpassage blocks; their questions carry group=<passage id>.
\documentclass[10pt]{article}
\usepackage[a4paper,margin=1.8cm]{geometry}
\usepackage[T1]{fontenc}
\usepackage{amsmath,amssymb,mathtools}
\usepackage{graphicx}
\usepackage[export]{adjustbox}
\usepackage{xcolor,needspace,keyval}
\usepackage{enumitem}
\usepackage{booktabs,array,tabularx}
\usepackage[hidelinks]{hyperref}
\setlength{\parindent}{0pt}
\setlength{\parskip}{3pt}
\makeatletter
\define@key{bq}{id}{\def\bq@id{#1}}\define@key{bq}{source}{\def\bq@src{#1}}\define@key{bq}{paper}{\def\bq@paper{#1}}
\define@key{bq}{qno}{\def\bq@qno{#1}}\define@key{bq}{type}{\def\bq@type{#1}}\define@key{bq}{subtopic}{\def\bq@sub{#1}}
\define@key{bq}{answer}{\def\bq@ans{#1}}\define@key{bq}{answer_src}{\def\bq@asrc{#1}}\define@key{bq}{working}{\def\bq@work{#1}}
\define@key{bq}{group}{\def\bq@group{#1}}\define@key{bq}{also_in}{\def\bq@also{#1}}
\define@key{bq}{difficulty}{\def\bq@diff{#1}}\define@key{bq}{diff_parts}{\def\bq@dparts{#1}}\define@key{bq}{diff_note}{\def\bq@dnote{#1}}
\newenvironment{bankq}[1][]{\def\bq@work{}\def\bq@also{}\def\bq@group{}\def\bq@diff{?}\def\bq@dparts{}\def\bq@dnote{}\setkeys{bq}{#1}%
  \par\needspace{4\baselineskip}\medskip\noindent{\small\textbf{\bq@id}\hfill\color{gray}[\bq@src, Paper \bq@paper, Q\bq@qno]
  \bq@type\ $\cdot$ \bq@sub\ $\cdot$ \textbf{D\bq@diff}}\par\nopagebreak\relax}
  {\par{\small\color{teal}Answer: \textbf{\bq@ans}\ifx\bq@work\@empty\else\ (solved: \bq@work)\fi
  \ifx\bq@also\@empty\else\par\color{gray}Also in: \bq@also\fi}\par}
\define@key{bp}{id}{\def\bp@id{#1}}\define@key{bp}{source}{\def\bp@src{#1}}\define@key{bp}{questions}{\def\bp@qs{#1}}
\newenvironment{bankpassage}[1][]{\setkeys{bp}{#1}\par\needspace{4\baselineskip}\bigskip\noindent
  {\small\textbf{Passage \bp@id}\hfill\color{gray}[\bp@src; questions \bp@qs]}\par\nopagebreak\relax\itshape}{\par}
\makeatother
\newcommand{\bfig}[2]{\par\smallskip{\centering\includegraphics[width=#1pt,max width=0.62\linewidth,max height=0.32\textheight]{img/#2}\par}\smallskip}
\newcommand{\bfigi}[2]{\includegraphics[width=#1pt,max width=\linewidth,max height=0.32\textheight,valign=c]{img/#2}}
\newcommand{\bmissing}[1]{\textit{[figure missing in the source conversion: \detokenize{#1}]}}
\newenvironment{qoptions}{\begin{itemize}[leftmargin=1.8em,itemsep=1pt,topsep=2pt,label={}]}{\end{itemize}}
\newcommand{\opt}[2]{\item[(#1)] #2}
\providecommand{\fwstart}{}\providecommand{\fwend}{}\providecommand{\qfullsplit}{}
\newcolumntype{L}[1]{>{\raggedright\arraybackslash}p{#1}}
\newcolumntype{C}[1]{>{\centering\arraybackslash}p{#1}}
\newcolumntype{R}[1]{>{\raggedleft\arraybackslash}p{#1}}
\providecommand{\sectionhdr}[1]{}\providecommand{\subsectionhdr}[1]{}


\begin{document}
{\Large\bfseries Sequences \& Series + Means --- Question Bank}\par\smallskip
{\small Mathematics chapter M03. Questions from all converted test papers, grouped by source series. Answers: printed key where it matches the paper, otherwise solved (answer\_src=solved, with working). Exact repeats across papers appear once, with also\_in listing the other copies.}\par\medskip\hrule

\section*{Score-I papers}
\begin{bankq}[id=SI12-P2-MATH-S3-Q07, source={Score-I - Test 12 (09 Mar 2022)}, paper=2, qno=7, type=matching, subtopic={GP}, answer={P$\to$3, Q$\to$2, R$\to$2}, answer_src=key, difficulty=7, diff_parts={K3 S2 I1 C1 F1}, diff_note={Maximum of f via u=5^x substitution gives a; integral limit gives r=1/3; S=a/(1-r)}]
Let the first term `$a$' of an infinite G.P. is the value of $x$ where function $f(x)=7+2x\log_e25-5^{x-1}-5^{2-x}$ has greatest value and common ratio `$r$' is equal to $\displaystyle\lim_{x\to0}\int_0^x\dfrac{t^2\,dt}{x^2\tan x}$. If $S$ is sum of infinite terms then
\begin{center}\footnotesize
\begin{tabular}{|p{0.6cm}|p{1.2cm}|p{0.6cm}|p{1.0cm}|}
\hline
 & \centering List-I & & \centering List-II\tabularnewline
\hline
(P) & $a$ & (1) & 4\\
(Q) & $\dfrac1r$ & (2) & 3\\
(R) & $S$ & (3) & 2\\
\hline
\end{tabular}
\end{center}
\end{bankq}

\section*{TOAS papers}
\begin{bankq}[id=TOAS01-P2-MATH-S1-Q04, source={TOAS - Test 01 (30 Jan 2022)}, paper=2, qno=4, type=single, subtopic={AM-GM-HM}, answer={C}, answer_src=key, difficulty=4, diff_parts={K1 S3 I1 C0 F0}, diff_note={AM-GM on each factor: 6 x 5 x 1}]
If $x,y,z$ are non-zero real numbers then the minimum value of expression
\[\left(\dfrac{(x^8+4x^4+1)(y^4+3y^2+1)(z^2+2z+2)}{x^4y^2}\right)\]
equals -
\begin{qoptions}
\opt{A}{12}
\opt{B}{24}
\opt{C}{30}
\opt{D}{48}
\end{qoptions}
\end{bankq}
\begin{bankq}[id=TOAS02-P1-MATH-S3-Q02, source={TOAS - Test 02 (13 Feb 2022)}, paper=1, qno=2, type=numerical, subtopic={Special series}, answer={1.5}, answer_src=solved, working={$T_n=(4^n-2^n)/2$, $S_n=(2\cdot4^n-3\cdot2^n+1)/3$ so $a=2,b=-3$, $|b/a|=3/2$}, difficulty=7, diff_parts={K2 S3 I1 C1 F1}, diff_note={Solve T_{n+1}=2T_n+4^n, sum the series, compare coefficients}]
Let $T_n$ be the $n^{\mathrm{th}}$ term of a series for $n\ge1$ and $T_1=1$, $T_{n+1}=2T_n+4^n$, $n\ge1$. If $S_n=$ Sum of first $n$ terms of the series and $S_n=\dfrac{a.4^n+b.2^n+1}{3}$ where $a,b\in I$, then the value of $\left|\dfrac ba\right|$ is
\end{bankq}
\begin{bankq}[id=TOAS03-P1-MATH-S1-Q05, source={TOAS - Test 03 (27 Feb 2022)}, paper=1, qno=5, type=single, subtopic={Special series}, answer={B}, answer_src=key, difficulty=2, diff_parts={K1 S2 I0 C0 F0}, diff_note={Odd terms = k - k/32}]
If $\dfrac1{1^5}+\dfrac1{2^5}+\dfrac1{3^5}+\dfrac1{4^5}+\ldots$ upto infinite terms $=k$, then value of $\left(\dfrac1{1^5}+\dfrac1{3^5}+\dfrac1{5^5}+\dfrac1{7^5}+\ldots\right.$ upto infinite terms$\left.\right)$ will be-
\begin{qoptions}
\opt{A}{$\dfrac k{32}$}
\opt{B}{$\dfrac{31k}{32}$}
\opt{C}{$\dfrac k{16}$}
\opt{D}{$\dfrac{15k}{16}$}
\end{qoptions}
\end{bankq}
\begin{bankq}[id=TOAS03-P1-MATH-S1-Q06, source={TOAS - Test 03 (27 Feb 2022)}, paper=1, qno=6, type=single, subtopic={AP}, answer={D}, answer_src=key, difficulty=2, diff_parts={K1 S1 I0 C1 F0}, diff_note={T6=S6-S5}]
If sum of first $n$ terms of a series is $S_n=3n^3+2n^2$, then $6^{\mathrm{th}}$ term is-
\begin{qoptions}
\opt{A}{309}
\opt{B}{720}
\opt{C}{101}
\opt{D}{295}
\end{qoptions}
\end{bankq}
\begin{bankq}[id=TOAS03-P1-MATH-S2-Q07, source={TOAS - Test 03 (27 Feb 2022)}, paper=1, qno=7, type=multi, subtopic={Special series}, answer={A,C,D}, answer_src=key, difficulty=7, diff_parts={K2 S3 I1 C1 F1}, diff_note={Split into geometric series plus sum of (r^2+1) and compare}]
Let $S_n=(2+2)+(2^2+5)+(2^3+10)+(2^4+17)+\ldots$ upto $n$ brackets. If $S_n=2^{n+A}+Bn^3+Cn^2+Dn+E\ \forall\ n\in\mathbb N$ (where $A,B,C,D,E$ are constants), then which of the following is/are true ?
\begin{qoptions}
\opt{A}{$A=1$}
\opt{B}{$B=\dfrac16$}
\opt{C}{$D=\dfrac76$}
\opt{D}{$E=-2$}
\end{qoptions}
\end{bankq}
\begin{bankq}[id=TOAS03-P1-MATH-S2-Q08, source={TOAS - Test 03 (27 Feb 2022)}, paper=1, qno=8, type=multi, subtopic={AM-GM-HM}, answer={C,D}, answer_src=key, difficulty=4, diff_parts={K1 S2 I1 C0 F1}, diff_note={AM-GM: x^2+27/x^2 >= 6 sqrt3 so values from 6 sqrt3+1 upward}]
If $x$ is a non-zero real, then value of $x^2+\dfrac{27}{x^2}+1$ can be-
\begin{qoptions}
\opt{A}{8}
\opt{B}{9}
\opt{C}{13}
\opt{D}{20}
\end{qoptions}
\end{bankq}
\begin{bankq}[id=TOAS03-P1-MATH-S3-Q05, source={TOAS - Test 03 (27 Feb 2022)}, paper=1, qno=5, type=numerical, subtopic={HP}, answer={2.00}, answer_src=key, difficulty=6, diff_parts={K2 S3 I1 C0 F1}, diff_note={HP gives AP of cosines; sum-to-product then 2 cos2t cos^2 a=1}]
If $\sec(2\theta-2\alpha),\ \dfrac12\mathrm{cosec}^2\theta,\ \sec(2\theta+2\alpha)$ are in H.P., then the value of $\sec2\theta\sec^2\alpha$ is
\end{bankq}
\begin{bankq}[id=TOAS03-P1-MATH-S3-Q06, source={TOAS - Test 03 (27 Feb 2022)}, paper=1, qno=6, type=numerical, subtopic={GP}, answer={2.25}, answer_src=solved, working={$a=b-d$, $c=b+d$; $(2+b)^2=(1+a)(1+c)$ gives $ac=(b+1)^2+2$ and $d^2=-2b-3\ge0$ so $b\le-\frac32$; minimum at $b=-\frac32$ ($a=b=c$): $ac=\frac94=2.25$ (the printed key 2.00 ignores $b\le-\frac32$)}, difficulty=7, diff_parts={K2 S3 I2 C0 F1}, diff_note={AP and GP conditions give ac=(b+1)^2+2 with b<=-3/2 from real a, c}]
Let $a,b,c$ are real number forming an A.P. and $(1+a),(2+b),(1+c)$ are in G.P., then minimum value of $ac$ is
\end{bankq}
\begin{bankq}[id=TOAS03-P2-MATH-S1-Q04, source={TOAS - Test 03 (27 Feb 2022)}, paper=2, qno=4, type=multi, subtopic={AM-GM-HM}, answer={A,B,C}, answer_src=key, difficulty=6, diff_parts={K2 S3 I1 C0 F1}, diff_note={AP of reciprocals and GP give b=-(a+c) and b^2=-2ac, test four relations}]
Let $a,b,c$ are non-zero real numbers such that $\dfrac1a,\dfrac1b,\dfrac1c$ are in arithmetic progression and $a,b,-2c$ are in geometric progression, then which of the following statement(s) is (are) must be true ?
\begin{qoptions}
\opt{A}{$a^2,b^2,4c^2$ are in geometric progression}
\opt{B}{$-2a,b,-2c$ are in arithmetic progression}
\opt{C}{$a^3+b^3+c^3-3abc=0$}
\opt{D}{$a^2,b^2,c^2$ are in harmonic progression}
\end{qoptions}
\end{bankq}
\begin{bankq}[id=TOAS03-P2-MATH-S1-Q06, source={TOAS - Test 03 (27 Feb 2022)}, paper=2, qno=6, type=multi, subtopic={AP}, answer={A,B,C}, answer_src=key, difficulty=4, diff_parts={K1 S3 I0 C0 F1}, diff_note={T_n=105-5n; check each pair and partial sums}]
Consider the series $S_n=100+95+90+\ldots+T_n$, then correct statement/s is/are (where $T_n$ represents $n^{\mathrm{th}}$ term) -
\begin{qoptions}
\opt{A}{$T_{16}+T_{26}=0$}
\opt{B}{$T_{11}+T_{31}=0$}
\opt{C}{$S_{30}-S_{11}=0$}
\opt{D}{$S_{15}-S_{27}=0$}
\end{qoptions}
\end{bankq}
\begin{bankq}[id=TOAS03-P2-MATH-S2-Q05, source={TOAS - Test 03 (27 Feb 2022)}, paper=2, qno=5, type=numerical, subtopic={GP}, answer={6.00}, answer_src=key, difficulty=5, diff_parts={K1 S3 I1 C0 F1}, diff_note={Sum and sum of squares give a(1-r+r^2)=9, r=1/2, S=24}]
Sum of first three terms of G. P. with non-integer common ratio is 21 and sum of their squares is 189. If $S$ denotes sum of infinite term of this series, then $\dfrac S4$ is equal to
\end{bankq}
\begin{bankq}[id=TOAS03-P2-MATH-S2-Q06, source={TOAS - Test 03 (27 Feb 2022)}, paper=2, qno=6, type=numerical, subtopic={AM-GM-HM}, answer={7.00}, answer_src=key, difficulty=6, diff_parts={K1 S3 I1 C1 F1}, diff_note={Expression equals 2 sin^2 x + 3/sin^2 x + a^2+1/a^2: minimise at sin^2x=1}]
Minimum possible integral value of the expression $\dfrac{2a^2\sin^4x+(a^4+1)\sin^2x+3a^2}{a^2\sin^2x}$ (wherever defined) is
\end{bankq}
\begin{bankq}[id=TOAS03-P2-MATH-S3-Q05, source={TOAS - Test 03 (27 Feb 2022)}, paper=2, qno=5, type=numerical, subtopic={Special series}, answer={7}, answer_src=key, difficulty=6, diff_parts={K1 S3 I1 C1 F1}, diff_note={Digit-cube sequence enters a 3-cycle; 102nd term found by position in cycle}]
A sequence start with 7120 as its first term and each subsequent term is the sum of cubes of the digits in the previous number. If $N$ is $102^{\mathrm{th}}$ term of this sequence, then sum of digits of $N$ is
\end{bankq}
\begin{bankq}[id=TOAS03-P2-MATH-S3-Q06, source={TOAS - Test 03 (27 Feb 2022)}, paper=2, qno=6, type=numerical, subtopic={Special series}, answer={1}, answer_src=key, difficulty=6, diff_parts={K1 S3 I2 C0 F1}, diff_note={1/(a_{n+1}-1)=1/(a_n-1)-1/a_n telescopes to 1}]
Let $a_1=2,\ a_{n+1}=a_n^2-a_n+1$ for $n\geq1$. Find $\displaystyle\sum_{n=1}^{\infty}\dfrac1{a_n}$
\end{bankq}
\begin{bankpassage}[id=TOAS04-P2-MATH-S2-G9-10, source={TOAS - Test 04 (15 Mar 2022)}, questions={9-10}]
The sequence $\{a_n\}_{n=1}^{+\infty}$ is defined by $a_1=0$ and $a_{n+1}=a_n+4n+3$, $n\ge1$
\end{bankpassage}
\begin{bankq}[id=TOAS04-P2-MATH-S2-Q09, source={TOAS - Test 04 (15 Mar 2022)}, paper=2, qno=9, type=single, subtopic={Special series}, answer={D}, answer_src=key, difficulty=3, diff_parts={K1 S3 I0 C0 F0}, diff_note={a_n=(n-1)(2n+3), a_50 mod 50}, group=TOAS04-P2-MATH-S2-G9-10]
The remainder when $a_{50}$ is divided by 50, is
\begin{qoptions}
\opt{A}{0}
\opt{B}{3}
\opt{C}{29}
\opt{D}{47}
\end{qoptions}
\end{bankq}
\begin{bankq}[id=TOAS04-P2-MATH-S2-Q10, source={TOAS - Test 04 (15 Mar 2022)}, paper=2, qno=10, type=single, subtopic={Special series}, answer={B}, answer_src=key, difficulty=6, diff_parts={K2 S3 I1 C1 F0}, diff_note={sqrt(a_n) ~ sqrt2 n so sums of 4^k vs 2^k terms give 683}, group=TOAS04-P2-MATH-S2-G9-10]
The value of
\[\lim_{n\to\infty}\frac{\sqrt{a_n}+\sqrt{a_{4n}}+\sqrt{a_{4^2n}}+\sqrt{a_{4^3n}}+\cdots+\sqrt{a_{4^{10}n}}}{\sqrt{a_n}+\sqrt{a_{2n}}+\sqrt{a_{2^2n}}+\sqrt{a_{2^3n}}+\cdots+\sqrt{a_{2^{10}n}}}\]
equals
\begin{qoptions}
\opt{A}{685}
\opt{B}{683}
\opt{C}{684}
\opt{D}{680}
\end{qoptions}
\end{bankq}
\end{document}
